\documentclass[11pt,a4paper]{article}
\usepackage[margin=26mm,headheight=15pt]{geometry}
\usepackage[T1]{fontenc}
\usepackage{amsmath,amsthm,mathtools}
\usepackage{newtxtext,newtxmath}
\usepackage{microtype}
\usepackage{booktabs,longtable,array}
\usepackage{xcolor}
\usepackage{listings}
\usepackage{enumitem}
\usepackage{needspace}
\usepackage{fancyhdr}
\usepackage{titlesec}
\usepackage{xurl}
\usepackage{hyperref}
\definecolor{accent}{HTML}{234765}
\definecolor{light}{HTML}{F2F5F7}
\hypersetup{colorlinks=true,linkcolor=accent,citecolor=accent,urlcolor=accent,
 pdftitle={Positive Real Asymptotic Inversion: Power-Log Series and Exact Irrational Supports},
 pdfauthor={OpenAI, prepared for Vladimir Reshetnikov}}
\pagestyle{fancy}\fancyhf{}
\fancyhead[L]{\small Positive real asymptotic inversion}
\fancyhead[R]{\small Theory and implementation}
\fancyfoot[C]{\thepage}
\renewcommand{\headrulewidth}{0.3pt}
\titleformat{\section}{\Large\bfseries\color{accent}}{\thesection}{0.7em}{}
\titleformat{\subsection}{\large\bfseries}{\thesubsection}{0.7em}{}
\setlength{\parskip}{4pt}
\setlength{\parindent}{1em}
\setlength{\emergencystretch}{2em}
\setlist{itemsep=2pt,topsep=4pt}
\numberwithin{equation}{section}
\newtheorem{theorem}{Theorem}[section]
\newtheorem{lemma}[theorem]{Lemma}
\newtheorem{proposition}[theorem]{Proposition}
\newtheorem{corollary}[theorem]{Corollary}
\theoremstyle{definition}
\newtheorem{definition}[theorem]{Definition}
\newtheorem{remark}[theorem]{Remark}
\newcommand{\R}{\mathbb R}
\newcommand{\C}{\mathbb C}
\newcommand{\N}{\mathbb N}
\newcommand{\OO}{\mathcal O}
\newcommand{\dd}{\mathrm d}
\newcommand{\D}{\mathcal D}
\newcommand{\kk}{\boldsymbol k}
\newcommand{\deltaa}{\boldsymbol\delta}
\newcommand{\degree}{\operatorname{deg}}
\newcommand{\val}{\operatorname{val}}
\newcommand{\coeff}[2]{[#1]#2}
\newcommand{\fall}[2]{#1^{\underline{#2}}}
\newcommand{\code}[1]{\texttt{#1}}
\lstdefinelanguage{Wolfram}{
 morekeywords={Module,With,If,Do,Table,Total,Sum,Return,Expand,FullSimplify,
  D,Log,Series,Normal,Get,Clear,Assumptions,True,False,Automatic,
  RealInverseExpansion,InverseResidual,LagrangeInverseTruncation},
 sensitive=true,morecomment=[s]{(*}{*)},morestring=[b]"}
\lstset{language=Wolfram,basicstyle=\ttfamily\small,keywordstyle=\color{accent},
 commentstyle=\itshape\color{black!60},stringstyle=\color{black!80},
 backgroundcolor=\color{light},frame=single,rulecolor=\color{black!15},
 breaklines=true,columns=fullflexible,keepspaces=true,showstringspaces=false,
 aboveskip=9pt,belowskip=9pt,xleftmargin=3pt,xrightmargin=3pt}

\begin{document}
\begin{titlepage}
\thispagestyle{empty}
\noindent\textcolor{accent}{\sffamily RESEARCH AND IMPLEMENTATION\hfill 7 SEPTEMBER 2026}\par
\vspace{18mm}
{\raggedright\huge\bfseries Positive real\par asymptotic inversion\par}
\vspace{5mm}
{\raggedright\Large Power--logarithmic series, exact irrational supports,\par and a Wolfram Language package\par}
\vspace{11mm}
\noindent Prepared for Vladimir Reshetnikov\par
\noindent OpenAI \quad\textbullet\quad Version 1.0.0
\vspace{10mm}
\begin{abstract}
The real inverse of $x+x^2(1+\log x)$ at $0^+$ is not a Taylor series,
and the inverse of $x+x^{\sqrt2}$ does not have a rational Puiseux support.
Both nevertheless admit a systematic, branch-explicit inversion calculus.
We develop it for every exact finite model
\[
 f(x)=a x^p+\sum_{j=1}^m x^{p+\eta_j}P_j(\log x),
 \qquad a,p,\eta_j>0,\quad P_j\in\R[L].
\]
After the positive substitution $v=x^p$, a local Lagrange--B\"urmann
formula gives every inverse coefficient as a finite product of first-order
operators on logarithmic polynomials. The output powers lie on a finitely
generated positive real semigroup, so irrational exponents and resonances
can be handled exactly. We prove branch existence and uniqueness, convergence
of the expansion for sufficiently small positive arguments, and a
next-weight remainder theorem that retains necessary powers of $\log y$.
An independent composition procedure checks the computed coefficients,
with the correct residual-to-inverse error conversion for $p\ne1$.
The accompanying package implements this finite-model theory, exposes
explicit remainder scales, and supplies a more general derivative-form
helper without pretending to certify its analytic hypotheses.
\end{abstract}
\vfill
\noindent\begin{tabular}{@{}p{0.20\textwidth}p{0.73\textwidth}@{}}
\toprule
Deliverables & This article and its editable source; commented package;
worked examples; native Wolfram tests; independent verification scripts.\\[4pt]
Verification & 66 independent mathematical and static checks passed,
including numerical comparisons at 180 decimal digits. The supplied
native Wolfram tests were not executed: the available service returned
HTTP 404 and no local Wolfram kernel was installed.\\
\bottomrule
\end{tabular}
\end{titlepage}
\tableofcontents
\clearpage

\section{The question and its concrete answers}\label{sec:answers}

The archived question asks for the positive real inverse near zero of
\begin{equation}\label{eq:question1}
 f(x)=x+x^2(1+\log x),\qquad x>0,
\end{equation}
and subsequently asks for the analogous expansion of
\begin{equation}\label{eq:question2}
 f_\alpha(x)=x+x^\alpha,\qquad \alpha=\sqrt2.
\end{equation}
The archive contains the question in Markdown and PDF; the Markdown
preserves both examples and their mathematical notation. The public
question is the corresponding primary source \cite{question}.
Its displayed computer outputs are historical reports by the questioner,
not results reproduced in a current Wolfram kernel for this article.

Put $L=\log y$. The first answer, through four correction orders, is
\begin{align}
 f^{-1}(y)={}&y-y^2(1+L)+y^3(2L^2+5L+3)\notag\\
 &-y^4\left(5L^3+\frac{41}{2}L^2+27L+\frac{23}{2}\right)\notag\\
 &+y^5\left(14L^4+\frac{241}{3}L^3+\frac{335}{2}L^2
                 +151L+\frac{299}{6}\right)
 +\OO\!\left(y^6(1+|L|)^5\right).
 \label{eq:log-answer}
\end{align}
In particular, the correct two-block statement is
\begin{equation}\label{eq:two-block}
 f^{-1}(y)=y-y^2(1+\log y)
          +\OO\!\left(y^3(1+|\log y|)^2\right).
\end{equation}
The error is \emph{not} $\OO(y^3)$: its next coefficient is a
nonconstant quadratic polynomial in $\log y$. More sharply,
\[
 \frac{f^{-1}(y)-y+y^2(1+\log y)}{y^3(\log y)^2}\longrightarrow2.
\]
This distinction matters even when a computer algebra system displays an
informal series-order marker that appears to suppress logarithms.

The second answer is exactly the expansion anticipated in the question:
\begin{align}
 f_{\sqrt2}^{-1}(y)={}&y-y^{\sqrt2}
       +\sqrt2\,y^{2\sqrt2-1}
       -\frac{6-\sqrt2}{2}\,y^{3\sqrt2-2}
       +\OO\!\left(y^{4\sqrt2-3}\right).
 \label{eq:irr-answer}
\end{align}
The next, explicitly computable term is
\[
 \frac{17\sqrt2-12}{3}\,y^{4\sqrt2-3}.
\]
All powers in these statements are evaluated on the positive real axis.
There is no need to choose an unspecified complex inverse and subsequently
try to repair its branch by assumptions on the output variable.

\Needspace{19\baselineskip}
\subsection{Immediate package usage}
After extracting the archive, use a path to the package file:
\begin{lstlisting}
Get["Kernel/RealInverseAsymptotics.wl"];
Clear[x, y, z];

r = RealInverseExpansion[
  x + x^2 (1 + Log[x]), {x, y}, 6];
r["Expression"]
r["RemainderScale"]
InverseResidual[r]["VanishingBelowCutoff"]

s = RealInverseExpansion[
  x + x^Sqrt[2], {x, z}, 4 Sqrt[2] - 3];
s["Expression"]
s["RemainderScale"]
\end{lstlisting}
The cutoff is \emph{exclusive}: the first call keeps output powers strictly
below $6$, and the second keeps exactly the four blocks displayed in
\eqref{eq:irr-answer}. The remainder field is a scale for a mathematical
big-$O$ statement, not a measured numerical error and not a substituted
\code{SeriesData} object. The residual call is designed to return
\code{True}; native-kernel execution status is recorded separately in
Section~\ref{sec:verification}.

\subsection{What is classical, and what is assembled here}
The coefficient identity used below is classical Lagrange inversion;
see the NIST DLMF \cite{dlmf} and Gessel's survey, especially
Theorem~2.1.1 \cite{gessel}. We do not claim to have invented that identity.
The contribution here is its complete adaptation to the specific
nonanalytic endpoint problem: a real branch condition, a positive-core
normalization, an exact irrational support calculus, a finite coefficient
algorithm, a proved logarithm-sensitive remainder, and an implementation
with an independent residual calculation.

The supplied repository's group README explicitly distinguishes the
near-identity operator form of Lagrange--B\"urmann from the classical
coefficient form and cautions against identifying results merely because
they share terminology \cite{repo}. The derivation in
Section~\ref{sec:lagrange} makes the bridge explicit. The repository's
canonical-volume README also separates exact, partial, and absent
formalization claims \cite{repovol}. No Lean verification of the present
package or of this article is asserted. Repository context consulted here
consists of those two READMEs, not a line-by-line review of the complete
consolidated volumes.

\section{The model, scales, and finite supports}\label{sec:scale}

\subsection{Exact input class}
Our main class is
\begin{equation}\label{eq:model}
 f(x)=a x^p+\sum_{j=1}^m x^{p+\eta_j}P_j(\log x),
 \quad a>0,\ p>0,\ \eta_j>0,\quad P_j\in\R[L].
\end{equation}
We first combine equal input powers and remove zero polynomials.
The sum is finite; the $\eta_j$ need not be integers, rational numbers,
or rationally independent. Coefficients of the $P_j$ can have either sign.
The leading coefficient $a$ is constant, rather than a polynomial in
$\log x$. A pure monomial, with $m=0$, is included and has the exact inverse
$(y/a)^{1/p}$. Statements below involving a minimum correction weight
or a first excluded weight assume $m\ge1$; they are unnecessary in
the pure-monomial case.

The model is initially regarded as an \emph{exact function}. An input
asymptotic expansion with an unknown remainder is a different object;
Section~\ref{sec:input-error} explains how that remainder affects the
inverse. Omitting it silently would turn a correct calculation for the
model into an unjustified assertion about the original function.

Define the normalization
\begin{equation}\label{eq:normalization}
 t=\frac ya,\qquad v=x^p,\qquad q=\frac1p,\qquad
 \delta_j=\frac{\eta_j}{p},\qquad
 Q_j(L)=\frac1a P_j\!\left(\frac Lp\right).
\end{equation}
Then the inverse problem is
\begin{equation}\label{eq:normalized}
 t=v+\psi(v),\qquad
 \psi(v)=\sum_{j=1}^m v^{1+\delta_j}Q_j(\log v),
 \qquad x=v^q.
\end{equation}
The positive-real identities $\log(x^p)=p\log x$ and
$x=(x^p)^{1/p}$ justify this substitution. They are branch assertions,
not unrestricted rules for complex powers.

\subsection{Power--logarithmic blocks}
A block is an expression $t^\gamma P(\log t)$, with $P$ a polynomial.
It is convenient to use the positive logarithmic envelope
\begin{equation}
 \ell(t)=1+|\log t|,\qquad 0<t<1.
\end{equation}
For every $\varepsilon>0$ and every fixed nonnegative integer $d$,
\begin{equation}\label{eq:power-dominates-log}
 t^\varepsilon\ell(t)^d\longrightarrow0.
\end{equation}
Indeed, putting $s=-\log t$ gives $e^{-\varepsilon s}(1+s)^d\to0$.
Consequently, an increase in the power of $t$ dominates any fixed change
in logarithmic degree. Within a single power, the highest nonzero power
of $\log t$ dominates the lower ones in magnitude.

A truncation by powers should retain entire coefficient polynomials.
For example, $y^3(2L^2+5L+3)$ is one block, but three scalar monomials.
A requested number of ``terms'' is therefore ambiguous unless the grading
is specified. The package uses an exclusive real-power cutoff rather than
an ambiguous term count.

\begin{lemma}[Uniqueness of a leading block]\label{lem:uniqueness}
Suppose $\gamma_1<\cdots<\gamma_r$ and $P_1\ne0$.
The finite sum $\sum_i t^{\gamma_i}P_i(\log t)$ is asymptotic to
$t^{\gamma_1}P_1(\log t)$ as $t\to0^+$.
In particular, distinct powers cannot cancel each other's leading
asymptotic behavior.
\end{lemma}
\begin{proof}
A nonzero polynomial satisfies $P_1(L)\sim cL^d$ as $L\to-\infty$.
Divide the sum by the first block and apply
\eqref{eq:power-dominates-log} to every remaining ratio.
\end{proof}

\subsection{Why irrational exponents cause no finite-truncation problem}
Let
\begin{equation}\label{eq:semigroup}
 \Gamma=\left\{\kk\mathbin\cdot\deltaa:
         \kk=(k_1,\ldots,k_m)\in\N_0^m\right\}.
\end{equation}
The inverse support will be contained in $q+\Gamma$.
The \emph{group} generated by irrationally related real numbers can be
dense; that observation is irrelevant to the \emph{positive semigroup}
used here.

\begin{lemma}[Local finiteness and finite fibers]\label{lem:finite}
For every finite $B$, there are only finitely many $\kk\in\N_0^m$
with $\kk\cdot\deltaa\le B$. In particular, $\Gamma\cap[0,B]$ is finite,
and each exponent has only finitely many multi-index representations.
\end{lemma}
\begin{proof}
Writing $\delta_* =\min_j\delta_j>0$ gives
$|\kk|\delta_*\le\kk\cdot\deltaa\le B$.
Alternatively, each coordinate satisfies $k_j\le B/\delta_j$.
Either inequality puts the relevant integer vectors in a finite box.
\end{proof}

Different vectors can still give the same power. Such a collision is
called a resonance here. It must be resolved by \emph{adding} the
corresponding logarithmic polynomials. Treating vectors as distinct output
powers would be mathematically wrong. Exact comparisons are essential
near resonances: rounding an exponent is not an acceptable way to decide
whether a term belongs below a cutoff.

\section{Selecting the positive real branch}\label{sec:branch}

\begin{theorem}[Existence, uniqueness, and leading behavior]\label{thm:branch}
For every model \eqref{eq:model}, there is $\varepsilon>0$ such that
$f$ is strictly increasing and positive on $(0,\varepsilon)$ and maps
that interval onto $(0,f(\varepsilon))$. Its inverse $x(y)$ on this
interval is the unique positive branch satisfying
\begin{equation}\label{eq:branch}
 \frac{x(y)}{(y/a)^{1/p}}\longrightarrow1\qquad(y\to0^+).
\end{equation}
In the normalized coordinate, $v(t)/t\to1$.
\end{theorem}
\begin{proof}
Write $f(x)=a x^p(1+r(x))$, where
\[
 r(x)=\frac1a\sum_jx^{\eta_j}P_j(\log x).
\]
By \eqref{eq:power-dominates-log}, $r(x)=o(1)$ and
\[
 xr'(x)=\frac1a\sum_jx^{\eta_j}
       (\eta_jP_j+P_j')(\log x)=o(1).
\]
Therefore
\begin{equation}\label{eq:derivative-leading}
 f'(x)=apx^{p-1}\left(1+r(x)+\frac{xr'(x)}p\right)
       =apx^{p-1}(1+o(1))>0
\end{equation}
for sufficiently small positive $x$. Also $f(x)>0$ there and
$f(x)\to0$. Continuity and strict monotonicity give the inverse interval.
Substitute $x=x(y)$ in $f(x)/(a x^p)\to1$ to obtain
$x(y)^p/(y/a)\to1$. The positive $p$th root gives
\eqref{eq:branch}. Any other positive branch tending to zero eventually
lies in the same strictly monotone interval and is identical to it.
\end{proof}

The theorem does not assert global invertibility for every model.
Coefficients of the corrections may create critical points farther from
zero. The asymptotic condition chooses a local real germ. A global inverse
exists for both examples in the question, as will be checked explicitly.

\subsection{Why assumptions on a complex inverse can miss this branch}
The ordinary analytic inverse theorem is centered at a point where the
forward map is analytic and has nonzero derivative; the DLMF states those
hypotheses explicitly \cite{dlmf}. In \eqref{eq:question1}, $x=0$ is not
an analytic point of $\log x$. A Taylor expansion about a nonzero complex
solution of $f(x)=0$ solves a \emph{different} local inverse problem.

For example, on compatible logarithmic sheets the nonzero equation
$x(1+\log x)=-1$ leads to
\[
 x=\exp\bigl(-1+W_k(-e)\bigr).
\]
An expansion with such a nonzero constant term cannot satisfy
$x(y)\sim y$. The correct branch condition is therefore built into the
coordinate normalization, rather than entrusted to an unspecified
\code{InverseFunction} representation. The current documentation describes
\code{InverseFunction} and restrictions of its domain \cite{inversefunction};
this article makes no claim that the historical calls in the question
still behave identically in every current release.

\section{Lagrange inversion away from the singular endpoint}\label{sec:lagrange}

\subsection{A local auxiliary parameter}
Fix $t>0$, temporarily. The function $\psi$ in \eqref{eq:normalized} is
analytic on a sufficiently small complex disk around this \emph{positive}
point, after choosing the logarithm continued from its real value.
Introduce an auxiliary parameter $s$:
\begin{equation}\label{eq:homotopy}
 V=t-s\psi(V),\qquad V(t,0)=t.
\end{equation}
This equation is regular at $s=0$, even though $\psi$ is not analytic at
$t=0$. Analyticity in $s$ is the useful regularity.

\begin{lemma}[Derivative-form Lagrange--B\"urmann]\label{lem:lagrange}
For an observable $\Phi$ analytic near the fixed positive $t$,
\begin{equation}\label{eq:observable}
 \Phi(V(t,s))=\Phi(t)+
 \sum_{n\ge1}\frac{(-s)^n}{n!}
 \frac{\dd^{n-1}}{\dd t^{n-1}}
       \bigl(\Phi'(t)\psi(t)^n\bigr)
\end{equation}
for sufficiently small $|s|$.
\end{lemma}
\begin{proof}
Set $w=V-t$. Then $w=sR(w)$ with $R(w)=-\psi(t+w)$.
The classical coefficient form of Lagrange inversion gives
\[
 [s^n]\Phi(t+w)
 =\frac1n[z^{n-1}]\Phi'(t+z)R(z)^n,
 \qquad n\ge1.
\]
This is precisely the form in Gessel's Theorem~2.1.1
\cite{gessel}. Taylor coefficient extraction at $z=0$ turns the right
side into
\[
 \frac{(-1)^n}{n!}
 \left.\frac{\dd^{n-1}}{\dd z^{n-1}}
   \bigl(\Phi'(t+z)\psi(t+z)^n\bigr)\right|_{z=0}.
\]
Differentiation in the translated variable gives \eqref{eq:observable}.
The analytic inverse theorem for $w-sR(w)=0$ justifies the convergent
local identity in $s$. Notice that $t$ was a positive parameter throughout;
no Taylor theorem was applied at the logarithmic singularity.
\end{proof}

For $\Phi(v)=v$, the formula becomes
\begin{equation}\label{eq:near-identity}
 V(t,s)=t+\sum_{n\ge1}\frac{(-s)^n}{n!}
                  \left(\frac{\dd}{\dd t}\right)^{n-1}\psi(t)^n.
\end{equation}
For the original inverse we need $\Phi(v)=v^q$, which gives
\begin{equation}\label{eq:power-observable}
 V(t,s)^q=t^q+
 \sum_{n\ge1}\frac{q(-s)^n}{n!}
 \left(\frac{\dd}{\dd t}\right)^{n-1}
           \bigl(t^{q-1}\psi(t)^n\bigr).
\end{equation}
The factor $q$ and the shift $q-1$ are important. Computing an inverse for
$v$ and then treating it as the inverse for $x$ would lose both.

The local lemma alone does not yet permit setting $s=1$ while $t\to0$.
A uniform smallness argument is still needed. Section~\ref{sec:convergence}
provides it, rather than silently identifying a formal perturbation series
with an asymptotic expansion of the actual inverse.

\subsection{Differentiating a logarithmic block}
Let $\D=\dd/\dd L$ act on coefficient polynomials.

\begin{lemma}[Block differentiation]\label{lem:block}
For any real $\lambda$, polynomial $P$, and integer $r\ge0$,
\begin{equation}\label{eq:block}
 \left(\frac{\dd}{\dd t}\right)^r
       \bigl(t^\lambda P(\log t)\bigr)
 =t^{\lambda-r}
       \prod_{j=0}^{r-1}(\D+\lambda-j)P(L)\bigg|_{L=\log t}.
\end{equation}
An empty product is the identity operator.
\end{lemma}
\begin{proof}
The first derivative is
$t^{\lambda-1}(\lambda P+P')(\log t)$.
Apply this identity successively. The exponent decreases by one at each
step, while the constant-coefficient polynomial differential operators
commute.
\end{proof}

Differentiation with respect to $t$ therefore requires only polynomial
arithmetic in $L$. This is the central computational simplification:
logarithms are neither discarded nor treated as constants when the
independent variable is differentiated.

\section{The coefficient theorem}\label{sec:coefficients}

For $\kk\in\N_0^m\setminus\{0\}$ write
\begin{equation}\label{eq:multi-notation}
 n=|\kk|=\sum_jk_j,\qquad
 \beta=\kk\cdot\deltaa,\qquad
 \kk!=\prod_jk_j!,\qquad
 Q^{\kk}(L)=\prod_jQ_j(L)^{k_j}.
\end{equation}
Define
\begin{equation}\label{eq:main-coefficient}
 \boxed{\displaystyle
 B_{\kk}^{(q)}(L)=
 \frac{q(-1)^n}{\kk!}
 \prod_{r=1}^{n-1}(\D+\beta+q+r)\,Q^{\kk}(L).}
\end{equation}
For $n=1$ this says $B_{\boldsymbol e_j}^{(q)}=-qQ_j$.

\begin{theorem}[Complete finite-model coefficient formula]\label{thm:coefficients}
The positive inverse of \eqref{eq:model} has the expansion
\begin{equation}\label{eq:full-series}
 \boxed{\displaystyle
 x(y)=t^q+\sum_{\kk\ne0}t^{q+\kk\cdot\deltaa}
                B_{\kk}^{(q)}(\log t),\qquad t=y/a.}
\end{equation}
Each coefficient is a polynomial. Equal powers are combined by
\begin{equation}\label{eq:aggregate}
 C_\beta(L)=\sum_{\kk\cdot\deltaa=\beta}B_{\kk}^{(q)}(L).
\end{equation}
Both the aggregate at a fixed power and the set of terms below a fixed
power cutoff are finite. The series converges as a sum of blocks for
sufficiently small $t>0$, and has the asymptotic remainder described in
Theorem~\ref{thm:remainder}.
\end{theorem}
\begin{proof}[Coefficient calculation]
The multinomial theorem gives
\[
 \psi(t)^n=\sum_{|\kk|=n}\frac{n!}{\kk!}
           t^{n+\beta}Q^{\kk}(\log t).
\]
Insert this expression in \eqref{eq:power-observable}. The initial
exponent in the differentiated block is $q-1+n+\beta$.
Lemma~\ref{lem:block}, with $r=n-1$, leaves the power $q+\beta$ and
produces the constants $q+\beta+1,\ldots,q+\beta+n-1$.
The multinomial $n!$ cancels the Lagrange denominator $n!$.
This proves \eqref{eq:main-coefficient}; Lemma~\ref{lem:finite} proves
all stated finiteness claims. The analytic and remainder assertions are
proved in the next section.
\end{proof}

\subsection{Degree and leading-logarithm checks}
Let $d_j=\degree Q_j$, and let $c_j$ be its leading coefficient.
For a nonzero multi-index block,
\begin{align}
 \degree B_{\kk}^{(q)}&=\sum_jk_jd_j,\label{eq:degree}\\
 [L^{\sum_jk_jd_j}]B_{\kk}^{(q)}
 &=\frac{q(-1)^n}{\kk!}
   \prod_{r=1}^{n-1}(\beta+q+r)\prod_jc_j^{k_j}.
 \label{eq:leading-log}
\end{align}
Every constant $\beta+q+r$ is positive, so the highest-degree term
cannot disappear inside one nonzero multi-index block. It can disappear
when several resonant blocks are added. This supplies an exact internal
check and explains why a conservative remainder degree should be computed
before assuming any cancellation.

One may expand the operator product using elementary symmetric
polynomials:
\begin{equation}\label{eq:operator-symmetric}
 \prod_{r=1}^{n-1}(\D+c_r)P
 =\sum_{j=0}^{\min(n-1,\degree P)}
    e_{n-1-j}(c_1,\ldots,c_{n-1})P^{(j)}.
\end{equation}
The implementation uses repeated first-order applications instead.
That approach avoids constructing an extra family of symmetric-polynomial
expressions and works directly in the coefficient polynomial ring.

\subsection{An alternative recursive construction}
The normalized equation can also be written with $v=t(1+u)$:
\begin{equation}\label{eq:unit-equation}
 0=u+\sum_jt^{\delta_j}(1+u)^{1+\delta_j}
                  Q_j\bigl(L+\log(1+u)\bigr),\qquad L=\log t.
\end{equation}
The derivative of its left side with respect to $u$ has constant term
$1$. Consequently the coefficient at each new support weight is determined
linearly by previously determined coefficients. Reversion by coefficient
comparison is therefore triangular, not a large nonlinear solve.

Formula~\eqref{eq:main-coefficient} solves this triangular process in a
finite nonrecursive form. The verification script uses the recursive
unit-series identities as an independent check; the package's residual
routine uses a third viewpoint, direct composition of a finite candidate.

\section{Convergence and logarithm-sensitive remainders}\label{sec:convergence}

\subsection{A uniform contraction in the auxiliary parameters}
Assume $m\ge1$, and put
\begin{equation}\label{eq:epsilon}
 \delta_* =\min_j\delta_j,\qquad d_* =\max_jd_j,\qquad
 \epsilon(t)=t^{\delta_*}\ell(t)^{d_*}.
\end{equation}
Thus $\epsilon(t)\to0$. Constants below may depend on the fixed model,
but not on sufficiently small positive $t$.

\begin{proposition}[Convergent perturbation expansion]\label{prop:convergence}
The series \eqref{eq:full-series} converges absolutely as a sum of its
multi-index block values for sufficiently small $t>0$. Its truncation
by total perturbation degree satisfies, for each fixed $N\ge0$,
\begin{equation}\label{eq:degree-error}
 x(y)-t^q-\sum_{1\le|\kk|\le N}
 t^{q+\kk\cdot\deltaa}B_{\kk}^{(q)}(\log t)
 =\OO\!\left(t^q\epsilon(t)^{N+1}\right).
\end{equation}
\end{proposition}
\begin{proof}
Introduce a separate complex parameter $s_j$ for each correction in
\eqref{eq:unit-equation}. On the closed disk $|u|\le1/2$, the functions
$(1+u)^{1+\delta_j}$ and $\log(1+u)$ are analytic and bounded, as are
the derivatives needed below. Polynomial growth gives uniform bounds
\[
 \left|t^{\delta_j}(1+u)^{1+\delta_j}
 Q_j(L+\log(1+u))\right|\le C\epsilon(t)
\]
and the same type of bound for the derivative in $u$.
Choose a fixed $K$ large enough that, whenever
\[
 |s_j|\le R(t):=\frac1{K\epsilon(t)}\quad(1\le j\le m),
\]
the map obtained by negating the sum in \eqref{eq:unit-equation}
sends $|u|\le1/2$ into $|u|\le1/4$ and has derivative bounded by
$1/2$ there. It is a contraction. Its successive approximations converge
uniformly on smaller parameter polydisks, so the unique fixed point is
holomorphic in the $s_j$. The observable $(1+u)^q$ is bounded by a
model-dependent constant $M$ on the same domain.

Multivariate Cauchy estimates imply
\[
 \left|[\boldsymbol s^{\kk}](1+u)^q\right|
 \le M R(t)^{-|\kk|}=M(K\epsilon(t))^{|\kk|}.
\]
By the coefficient calculation, the left side is
$|t^{\kk\cdot\deltaa}B_{\kk}^{(q)}(L)|$.
There are $\binom{n+m-1}{m-1}$ vectors with $|\kk|=n$.
For $K\epsilon(t)\le1/2$, summing these bounds proves absolute
convergence at $s_1=\cdots=s_m=1$, and the tail after total degree
$N$ is bounded by a constant depending on $N$ times
$\epsilon(t)^{N+1}$. Multiplication by $t^q$ gives
\eqref{eq:degree-error}.

At real $s_j=1$, the fixed point is real and $1+u>0$. It solves the
original normalized equation and has $u\to0$, so
Theorem~\ref{thm:branch} identifies it with the required inverse.
\end{proof}

This is a convergence theorem for the \emph{exact finite model}, not for
arbitrary formal transseries. It does not assert a universal numerical
convergence threshold. In particular, the neighborhood can shrink when
parameters change, and the estimates are not uniform as
$\delta_*\to0^+$.

\subsection{The first excluded weight}
Let $R>q$ be an output-power cutoff and set $B=R-q$.
Define
\begin{align}
 \beta_+&=\min\{\beta\in\Gamma:\beta\ge B\},\label{eq:frontier}\\
 D_+&=\max\left\{\sum_jk_jd_j:\kk\cdot\deltaa=\beta_+\right\}.
 \label{eq:frontier-degree}
\end{align}
Both are finite and computable. The set defining $\beta_+$ is nonempty
because it contains an integer multiple of $\delta_*$. More precisely,
\begin{equation}\label{eq:frontier-search}
 \beta_+\le M_B:=\left\lceil\frac B{\delta_*}\right\rceil\delta_*
 <B+\delta_*
\end{equation}
for every $B>0$. When $B/\delta_*$ is an integer, $M_B=B$.
The implementation enumerates the closed weighted simplex up to $M_B$,
so it sees the boundary even though the output cutoff is exclusive.

\begin{theorem}[Next-weight remainder]\label{thm:remainder}
Let
\begin{equation}\label{eq:truncated}
 X_R(y)=t^q+\sum_{\kk\cdot\deltaa<R-q}
            t^{q+\kk\cdot\deltaa}B_{\kk}^{(q)}(\log t),
\end{equation}
where the zero vector is omitted from the sum. Then
\begin{equation}\label{eq:sharp-order}
 x(y)-X_R(y)=\OO\!\left(t^{q+\beta_+}\ell(t)^{D_+}\right).
\end{equation}
In fact,
\begin{equation}\label{eq:first-omitted}
 x(y)-X_R(y)=t^{q+\beta_+}C_{\beta_+}(\log t)
             +o\!\left(t^{q+\beta_+}\ell(t)^{D_+}\right).
\end{equation}
If the aggregate $C_{\beta_+}$ is nonzero, its actual degree and leading
coefficient determine the leading asymptotic error.
\end{theorem}
\begin{proof}
Choose a fixed integer $N$ with $(N+1)\delta_*>\beta_+$.
The tail in \eqref{eq:degree-error} is
\[
 \OO\!\left(t^{q+(N+1)\delta_*}\ell(t)^{(N+1)d_*}\right)
 =o\!\left(t^{q+\beta_+}\ell(t)^{D_+}\right)
\]
by the positive power gap. All indices with weight at most $\beta_+$
have degree at most $N$. Among the remaining finitely many indices of
degree at most $N$, those at $\beta_+$ give exactly the polynomial
$C_{\beta_+}$; every larger weight is negligible by
\eqref{eq:power-dominates-log}. This proves both statements. If $C_{\beta_+}\ne0$, the same positive
power-gap argument works with $\degree C_{\beta_+}$ in place of $D_+$,
so that polynomial also determines the leading asymptotic error.
\end{proof}

The proof is more than the observation that a first omitted exponent
exists. Logarithmic degrees grow with perturbation degree; an uncontrolled
infinite tail cannot be bounded merely by looking at its smallest power.
The finite-degree split and Proposition~\ref{prop:convergence} supply the
missing uniform control.

The package uses the safe degree $D_+$ before cancellations. It does not
scan indefinitely to find the first \emph{nonzero} omitted aggregate.
Thus its remainder may be conservative in resonant or specially tuned
examples, but it remains valid. When the first excluded power equals the
cutoff and $D_+>0$, replacing its scale by a bare $t^R$ is generally
incorrect.

\section{Worked examples and coefficient structure}\label{sec:examples}

\subsection{The logarithmic example: a global real inverse}
For $f(x)=x+x^2(1+\log x)$,
\begin{equation}
 f'(x)=1+x(3+2\log x),\qquad f''(x)=5+2\log x.
\end{equation}
The derivative $f'$ decreases up to $x=e^{-5/2}$ and increases thereafter.
Its global minimum on $(0,\infty)$ is
\begin{equation}\label{eq:global-slope}
 m_0=1-2e^{-5/2}>0.
\end{equation}
Since $f(0^+)=0$ and $f(x)\to\infty$, the function is a bijection of
$(0,\infty)$ onto itself. Thus the local branch constructed above is
also the unique global positive inverse.

Here $p=a=q=\delta_1=1$, $Q_1(L)=1+L$, and the expansion has the form
\begin{equation}\label{eq:log-polys}
 f^{-1}(y)=y+\sum_{n\ge1}y^{n+1}A_n(\log y),
 \qquad
 A_n(L)=\frac{(-1)^n}{n!}
       \prod_{j=2}^{n}(\D+n+j)(1+L)^n.
\end{equation}
For example,
\[
 A_2(L)=\tfrac12(\D+4)(1+L)^2
        =(1+L)+2(1+L)^2=3+5L+2L^2,
\]
and
\[
 A_3(L)=-\tfrac16(\D+5)(\D+6)(1+L)^3.
\]
These computations give \eqref{eq:log-answer}. For any fixed $N\ge0$,
\begin{equation}\label{eq:log-N-error}
 f^{-1}(y)=y+\sum_{n=1}^Ny^{n+1}A_n(\log y)
          +\OO\!\left(y^{N+2}\ell(y)^{N+1}\right).
\end{equation}
The next weight is exactly $N+1$, and its logarithmic degree is $N+1$.

\paragraph{Catalan leading coefficients.}
The constant factors in \eqref{eq:log-polys} range from $n+2$ to $2n$.
Therefore
\begin{equation}\label{eq:catalan}
 [L^n]A_n(L)=(-1)^n\frac{(2n)!}{n!(n+1)!}
            =(-1)^n\operatorname{Cat}_n.
\end{equation}
This identity is useful both conceptually and computationally. If one
freezes the slowly varying factor at its leading logarithmic value,
the dominant coefficients resemble inversion of a quadratic perturbation;
\eqref{eq:catalan} states the precise surviving highest-logarithm law.
The lower logarithmic powers are not obtained by freezing the logarithm:
they also contain the effect of its derivatives.

\paragraph{A triangular recurrence that checks the formula.}
Introduce a formal variable $s$ independent of $L$, and put
$U(s)=\sum_{n\ge1}A_n(L)s^n$. The unit equation reads
\begin{equation}\label{eq:log-recursion}
 U=-s(1+U)^2\bigl(1+L+\log(1+U)\bigr).
\end{equation}
It follows that
\[
 A_n(L)=-[s^{n-1}](1+U_{n-1})^2
               \bigl(1+L+\log(1+U_{n-1})\bigr),
\]
where $U_{n-1}=\sum_{j=1}^{n-1}A_j(L)s^j$.
The right side depends only on earlier coefficients. The independent
verification checks this recurrence against the differential-operator
formula through $n=8$.

\paragraph{Why this is not a Taylor series.}
The inverse extends continuously to zero and has right derivative $1$.
For $y>0$, ordinary differentiation of the inverse gives
\begin{equation}
 (f^{-1})''(y)=
 -\frac{5+2\log(f^{-1}(y))}
 {\bigl(1+f^{-1}(y)(3+2\log(f^{-1}(y)))\bigr)^3}
 \longrightarrow+\infty.
\end{equation}
Hence it does not have a finite second right derivative at zero, much
less an analytic Taylor expansion there. The singularity is compatible
with, and directly visible in, the $y^2\log y$ term.

\subsection{A single arbitrary real power}
Consider
\begin{equation}\label{eq:single-power}
 F_{\alpha,c}(x)=x+c x^\alpha,\qquad \alpha>1,\quad c\in\R.
\end{equation}
Its positive inverse exists near zero for every fixed real $c$.
The general formula reduces to
\begin{align}
 F_{\alpha,c}^{-1}(y)
 &=y+\sum_{n\ge1}(-c)^n b_n(\alpha)
                   y^{1+n(\alpha-1)},\label{eq:alpha-series}\\
 b_n(\alpha)
 &=\frac{\fall{(n\alpha)}{n-1}}{n!}
   =\frac1n\binom{n\alpha}{n-1}
   =\frac1{1+n(\alpha-1)}\binom{n\alpha}{n}.
 \label{eq:alpha-coef}
\end{align}
The falling factorial is a finite product, so the coefficient formula
needs no special-function evaluation and is meaningful for every real
$\alpha>1$.

The first coefficients are
\begin{align*}
 F_{\alpha,c}^{-1}(y)={}&y-cy^\alpha+c^2\alpha\,y^{2\alpha-1}
 -\frac{c^3\alpha(3\alpha-1)}2y^{3\alpha-2}\\
 &+\frac{c^4\alpha(4\alpha-1)(2\alpha-1)}3y^{4\alpha-3}
 +\OO\!\left(y^{5\alpha-4}\right).
\end{align*}
At $\alpha=\sqrt2$ and $c=1$, the cubic correction coefficient is
$-(6-\sqrt2)/2$, reproducing the question. The exponent increment
$\sqrt2-1$ is irrational, but positive; support enumeration is simply
enumeration of its nonnegative integer multiples.

For $c=1$, $F_{\alpha,1}'(x)=1+\alpha x^{\alpha-1}>0$ on the entire
positive axis. Thus the inverse in this case is globally unique. For
negative $c$, the construction remains local and makes no corresponding
global claim.

\paragraph{An exact convergence radius in the small parameter.}
Put $s=c y^{\alpha-1}$ and $F_{\alpha,c}^{-1}(y)=yW(s)$.
Then $W+sW^\alpha=1$, $W(0)=1$, and
\[
 W(s)=1+\sum_{n\ge1}(-1)^n b_n(\alpha)s^n.
\]
Its Taylor radius is
\begin{equation}\label{eq:alpha-radius}
 R_\alpha=\frac{(\alpha-1)^{\alpha-1}}{\alpha^\alpha}.
\end{equation}
To verify this, write
\[
 b_n(\alpha)=\frac{\Gamma(n\alpha+1)}
 {\Gamma(n+1)\Gamma(n(\alpha-1)+2)}.
\]
Stirling's formula gives
\begin{equation}\label{eq:alpha-coef-asymp}
 b_n(\alpha)\sim
 \frac{\sqrt\alpha}{\sqrt{2\pi}(\alpha-1)^{3/2}}
 R_\alpha^{-n}n^{-3/2}.
\end{equation}
The root test proves \eqref{eq:alpha-radius}. The same value is suggested
by the simultaneous equations $W+sW^\alpha=1$ and
$1+\alpha sW^{\alpha-1}=0$, which give
$W=\alpha/(\alpha-1)$ and $s=-R_\alpha$.
Thus \eqref{eq:alpha-series} is not merely formal: it converges whenever
$|c|y^{\alpha-1}<R_\alpha$. At $\alpha=2$ this specializes to the familiar
Catalan radius $1/4$.

\subsection{Mixed logarithmic and irrational generators}
Let
\[
 F(x)=x+b x^{\sqrt2}+c x^{3/2}(1+\log x),\qquad b,c\in\R.
\]
The two relative weights are $\delta_1=\sqrt2-1$ and
$\delta_2=1/2$. Below output power $2$, the inverse is
\begin{align}
 F^{-1}(y)={}&y-b y^{\sqrt2}-c y^{3/2}(1+L)
 +\sqrt2\,b^2y^{2\sqrt2-1}\notag\\
 &+bc\,y^{\sqrt2+1/2}
       \left[\left(\sqrt2+\frac32\right)(1+L)+1\right]
 +\OO\!\left(y^2\ell(y)^2\right).
 \label{eq:mixed-example}
\end{align}
The mixed coefficient corresponds to $\kk=(1,1)$ and is obtained by
one application of $\D+\sqrt2+3/2$ to $bc(1+L)$.
This example shows why neither ordinary integer order nor total
perturbation degree alone is the requested output ordering. A single
weighted support calculation handles both types of term.

\subsection{Resonances}
For
\[
 F(x)=x+A x^{3/2}+B x^2,
\]
the weight $1$ has two representations: $2(1/2)$ and $1$.
The inverse begins
\begin{align}
 F^{-1}(y)={}&y-Ay^{3/2}
 +\left(\frac32A^2-B\right)y^2\notag\\
 &+\left(-\frac{21}{8}A^3+\frac72AB\right)y^{5/2}
 +\OO(y^3).
 \label{eq:resonance-example}
\end{align}
The coefficient of $y^2$ is not a choice between two candidate answers:
it is their sum. It may vanish, for example when $B=3A^2/2$.
That cancellation does not alter the semigroup, but it changes the
nonzero support and can make the generic error envelope conservative.

\subsection{A non-unit leading power}
For $F(x)=x^2+x^3$, use $v=x^2$, $q=1/2$, and
$t=v+v^{3/2}$. Applying the observable formula, rather than returning
$v$ itself, gives
\begin{equation}\label{eq:power-core-example}
 F^{-1}(y)=y^{1/2}-\frac12y+\frac58y^{3/2}-y^2
             +\frac{231}{128}y^{5/2}+\OO(y^3).
\end{equation}
The derivative of the forward map tends to zero, so a residual and an
inverse error no longer have the same power. A candidate with inverse
error of order $y^3$ has forward residual of order $y^{7/2}$.
This is the reason for the cutoff conversion in the residual checker.

For a positive leading coefficient $a\ne1$, the natural coordinate is
$t=y/a$. The package deliberately leaves $(y/a)^\gamma$ and
$\log(y/a)$ in this form. It does not use the complex-unsafe blanket
transformation of powers often associated with \code{PowerExpand}.
Under $a>0$ and $y>0$, conversion to powers of $y$ and polynomials in
$\log y-\log a$ is mathematically legitimate, but not needed internally.

\section{Residual certificates and uncertain input expansions}\label{sec:residual}

\subsection{The correct transport of a residual}
Let $x(y)$ denote the exact inverse and let $X(y)>0$ be a candidate with
$X(y)/t^q\to1$. Define the exact normalized residual
\begin{equation}\label{eq:residual-definition}
 E(y)=\frac{f(X(y))}{a}-t.
\end{equation}

\begin{proposition}[Residual-to-inverse transport]\label{prop:transport}
For a fixed model \eqref{eq:model},
\begin{equation}\label{eq:transport}
 |X(y)-x(y)|\le C t^{q-1}|E(y)|
\end{equation}
for sufficiently small positive $y$. More precisely, at points where the
residual is nonzero,
\begin{equation}\label{eq:transport-first}
 X(y)-x(y)=\frac1p t^{q-1}E(y)(1+o(1)).
\end{equation}
Thus a residual of order $\OO(t^A\ell(t)^d)$ gives inverse error
$\OO(t^{A+q-1}\ell(t)^d)$.
\end{proposition}
\begin{proof}
The mean value theorem gives
$f(X)-f(x)=f'(\xi)(X-x)$ for a point $\xi$ between $X$ and $x$.
Both endpoints are asymptotic to $t^q$, hence so is $\xi$.
Equation~\eqref{eq:derivative-leading} gives
\[
 f'(\xi)=ap\,t^{q(p-1)}(1+o(1))
        =ap\,t^{1-q}(1+o(1)).
\]
Divide the residual identity by this derivative.
\end{proof}

An inverse jet retaining powers below $R$ should therefore have vanishing
forward residual below
\begin{equation}\label{eq:residual-cutoff}
 C=R+1-q.
\end{equation}
For $p=1$ this equals $R$, but for a square-root inverse it equals
$R+1/2$. Testing a residual at an unconverted cutoff can be too weak or
unnecessarily strong.

\subsection{Independent finite composition}
Write the candidate as $X=t^q(1+U)$, where $U$ is a finite power--log jet
with strictly positive least power $\nu$. An original block transforms as
\begin{equation}\label{eq:compose-block}
 X^\alpha P(\log X)
 =t^{q\alpha}(1+U)^\alpha
       P\bigl(qL+\log(1+U)\bigr).
\end{equation}
Introduce an ordinary auxiliary variable $z$ and expand the analytic unit
\[
 H(z,L)=(1+z)^\alpha P\bigl(qL+\log(1+z)\bigr)
\]
at $z=0$. The term containing $U^j$ has power at least $j\nu$.
Therefore only finitely many $j$ can contribute below a residual cutoff
$C$: it suffices to compute through
$\lfloor(C-q\alpha)/\nu\rfloor$, then discard the boundary if necessary.
Sparse multiplication of $U$ is truncated at each step. Because all
powers of $U$ are positive, an already discarded term can never return
to a lower power after another multiplication.

This composition procedure does not use the inverse coefficient formula.
It is the algorithm implemented by \code{InverseResidual}. It returns
a finite \emph{residual jet}, not the entire nonlinear expression
$f(X)-y$. In particular, a returned residual expression of zero means
``all retained coefficients vanish,'' not ``the finite approximation is
an exact inverse.'' The function checks the stored \code{"Blocks"}
field; editing only a separate \code{"Expression"} field does not change
what is checked.

\subsection{Numerical a posteriori bounds for the first example}
For \eqref{eq:question1}, the global lower bound
\eqref{eq:global-slope} is available. Every positive trial point $X$ obeys
\begin{equation}\label{eq:global-numeric-bound}
 |X-f^{-1}(y)|\le
 \frac{|X+X^2(1+\log X)-y|}{1-2e^{-5/2}}.
\end{equation}
If the exact residual is positive, the inverse lies between
$X-|f(X)-y|/m_0$ and $X$; if it is negative, it lies between $X$ and
$X+|f(X)-y|/m_0$. One may intersect these intervals with $(0,\infty)$.

Such a statement is a numerical certificate only when the residual is
itself rigorously enclosed, for example with appropriately controlled
interval arithmetic. A high-precision floating-point computation alone
is a diagnostic, not a certified enclosure. The delivered package does
not claim to manufacture interval certificates or explicit constants for
its asymptotic big-$O$ metadata.

More generally, suppose a derivative lower bound $m>0$ is verified on an
interval containing $X-\rho$ and $X+\rho$, and $|f(X)-y|\le m\rho$.
Integration of the derivative bound gives opposite weak signs at the two
endpoints. Continuity provides a root in the interval, and strict
monotonicity makes it unique. This separates a genuine root-existence
certificate from a merely small residual.

\subsection{Transporting an input remainder}\label{sec:input-error}
Suppose the actual function is $F=f+r$, where $f$ is the finite model,
\begin{equation}\label{eq:input-remainder}
 r(x)=\OO\!\left(x^\lambda\ell(x)^M\right),\qquad \lambda>p,
\end{equation}
and $F'(x)=apx^{p-1}(1+o(1))$. The latter condition follows, for example,
from $r'(x)=o(x^{p-1})$. Let $x_F$ and $x_f$ be the positive inverse
branches. Evaluating $F$ at $x_f$ and applying the same mean-value argument
gives
\begin{equation}\label{eq:input-transport}
 x_F(y)-x_f(y)=
 \OO\!\left(t^{q\lambda+q-1}\ell(t)^M\right).
\end{equation}
Thus input uncertainty and algorithmic truncation must both be included.
If the missing forward term is of order $x^{p+\eta}$, its possible effect
on the inverse begins at power $q+\eta/p$.

The derivative condition is not cosmetic. The function
\[
 F(x)=x+x^2\sin(x^{-3})
\]
satisfies $F(x)/x\to1$, but its derivative changes sign arbitrarily close
to zero. Mere relative closeness to a monotone leading term does not
supply a single monotone inverse germ. The finite power--log model avoids
this problem because its derivatives have controlled scales.

\paragraph{Example: an elementary function first reduced to a jet.}
Near zero,
\[
 x+\sin x=2x-\frac{x^3}{6}+\frac{x^5}{120}+\OO(x^7),
\]
with the corresponding derivative control. Inverting the displayed model
and transporting the input remainder gives
\[
 (x+\sin x)^{-1}(y)=\frac y2+\frac{y^3}{96}
                       +\frac{y^5}{1920}+\OO(y^7).
\]
The package rejects an unexpanded \code{Sin[x]} in its exact-model front
end, but accepts the displayed polynomial model. The theorem explains
precisely what one may conclude about the original elementary function;
applying \code{Normal} to a forward series without recording its error
would not explain this by itself.

\section{Algorithms and implementation contracts}\label{sec:algorithms}

\subsection{Inversion algorithm}
The principal routine implements the following finite calculation.
\begin{enumerate}[label=\arabic*.]
\item Parse and combine the exact input blocks. Check real coefficients,
positive leading coefficient, positive leading power, and exact ordered
numerical exponents. Normalize by \eqref{eq:normalization}.
\item Set $B=R-q$ and compute the frontier bound $M_B$ in
\eqref{eq:frontier-search}. Enumerate the integer vectors
$\kk\cdot\deltaa\le M_B$, pruning recursively by the remaining weight
budget. Impose explicit resource limits.
\item For vectors with $0<\kk\cdot\deltaa<B$, form $Q^{\kk}$ and
apply the $n-1$ first-order operators in \eqref{eq:main-coefficient}.
Combine equal exponents and sort them by exact comparison.
\item Find $\beta_+$ and $D_+$ among the excluded frontier vectors.
Return the expression, coefficient blocks, model data, branch convention,
and remainder scale from Theorem~\ref{thm:remainder}.
\end{enumerate}
All indices below the requested cutoff are included because the frontier
bound is at least $B$. No index beyond that bound can have a smaller
excluded weight. The coefficient theorem proves the formula applied to
each included vector; exact grouping proves resonance handling; the
remainder theorem proves the meaning of the returned error scale.
These observations establish the mathematical correctness of the
algorithm independently of any particular CAS evaluation.

\subsection{Representation and differentiation}
Internally, a block is a pair
\begin{lstlisting}
{exponent, polynomialInLogCoordinate}
\end{lstlisting}
where \code{LogCoordinate} is a protected public indeterminate.
A coefficient operator is implemented by repeated updates of the form
\begin{lstlisting}
poly = Expand[D[poly, LogCoordinate] +
              (beta + q + j) poly];
\end{lstlisting}
The numerical variable is not present in these polynomial calculations.
Only when constructing an output expression is \code{LogCoordinate}
replaced by \code{Log[y/a]} and the block multiplied by
\lstinline|(y/a)^exponent|.

The parser replaces only the literal logarithm \code{Log[x]} by a fresh
indeterminate, expands finite algebraic sums, and extracts factors that
are powers of the input variable. It deliberately does not rewrite
arbitrary nested logarithms or fractional powers of sums. Positivity of
$x$ makes many additional identities possible, but a controlled parser
should not obtain them through indiscriminate complex power rules.

\subsection{Exact order and failure behavior}
The comparison routine tries to prove $\alpha-\beta=0$, $<0$, or $>0$
using exact simplification under the supplied assumptions. If none can
be established, it returns a \code{Failure}. There is no decimal fallback
for support ordering. Likewise, integer bounds for enumeration must be
established as exact integers.

The implemented front end requires exponent constants such as integers,
rationals, radicals, or other exact real numerical constants that the
underlying CAS can compare. It does not accept unrestricted symbolic
exponents even when a partial inequality, such as $\alpha>1$, is known.
That restriction is an implementation boundary, not a restriction of the
mathematical theorem. Coefficients may be symbolic if their reality and
the leading coefficient's positivity follow from \code{Assumptions}.

Unsupported input, unproved sign conditions, an invalid cutoff, and
exceeded enumeration limits return a structured \code{Failure}; they do
not produce a guessed branch or a silently truncated coefficient set.
The package reserves \code{LogCoordinate} and requires distinct,
unassigned input and output symbols.

\subsection{Cost and current engineering limits}
Let $K$ be the number of enumerated vectors and
$N_{\max}=\lfloor M_B/\delta_*\rfloor$. A crude bound is
\begin{equation}
 K\le\prod_{j=1}^m
     \left(1+\left\lfloor\frac{M_B}{\delta_j}\right\rfloor\right).
\end{equation}
For fixed weights and fixed $m$, this is polynomial growth in the cutoff,
though the dimension $m$ matters substantially. A coefficient of total
degree $n$ requires $n-1$ first-order polynomial operators and has
logarithmic degree at most $nd_*$.

The reference implementation prioritizes transparency and exact
comparisons. It uses simple list-based aggregation, so exponent merging
can entail quadratically many comparisons in the number of blocks.
It is not an optimized high-order transseries engine. Large logarithmic
degrees can also cause substantial exact-expression growth. Defaults
limit total frontier order to $64$ and enumerated vectors to $100000$;
users can change these bounds explicitly. The limits are safety controls,
not performance promises.

Useful future optimizations include canonical algebraic-number keys,
cached polynomial powers, memoized operator products, direct aggregation
by a rational lattice when one exists, and Newton lifting in the same
sparse jet algebra. None is needed to make the present finite calculation
well-defined or to obtain the requested examples.

\section{Wolfram Language interface}\label{sec:interface}

\subsection{\texttt{RealInverseExpansion}}
The interface is
\begin{lstlisting}
RealInverseExpansion[expr, {x, y}, cutoff,
  Assumptions -> True,
  "MaxMultiIndices" -> 100000,
  "MaxTotalOrder" -> 64]
\end{lstlisting}
It accepts a finite sum of terms $x^\alpha P(\log x)$ with a positive
pure-power leading term. Inputs containing approximate real numbers are
rejected; construct the expansion exactly, then substitute numerical
parameters. The cutoff must be an exact real numerical constant strictly
greater than $q=1/p$.

The principal result fields are as follows.
\begin{longtable}{@{}p{0.34\textwidth}p{0.61\textwidth}@{}}
\toprule
Field & Meaning\\\midrule
\endfirsthead
\toprule Field & Meaning\\\midrule\endhead
\code{"Expression"} & The finite inverse approximation, evaluated in the
coordinate $t=y/a$.\\[3pt]
\code{"Blocks"} & Ordered pairs $\{q+\beta,C_\beta(L)\}$, using
\code{LogCoordinate} for $L$. These are the data checked by the residual
routine.\\[3pt]
\code{"RemainderScale"} & A scale $t^{q+\beta_+}\ell(t)^{D_+}$ such that
the error is big-$O$ of that scale for the exact model. It is zero for a
pure monomial.\\[3pt]
\code{"RemainderExponent"} & $q+\beta_+$, or \code{Infinity} for the exact
pure-monomial case.\\[3pt]
\code{"RemainderLogDegree"} & The safe degree $D_+$.\\[3pt]
\code{"Cutoff"} & The exclusive output-power cutoff.\\[3pt]
\code{"LeadingInversePower"} & $q=1/p$.\\[3pt]
\code{"ScaledVariable"} & The expression $y/a$.\\[3pt]
\code{"Model"} & Original blocks, leading power and coefficient,
normalized corrections, and the variable names.\\[3pt]
\code{"Assumptions"} & The coefficient and parameter assumptions used
for exact simplification. Positive $y$ is part of the branch contract.\\[3pt]
\code{"Branch"} & The positive-real convention
$x/(y/a)^{1/p}\to1$ as $y\to0^+$.\\[3pt]
\code{"Exact"} & \code{True} for the pure-monomial special case;
\code{False} otherwise. It does not indicate numerical precision.\\[3pt]
\code{"EnumeratedMultiIndices"} & Number of nonzero frontier vectors
examined.\\[3pt]
\code{"ContributingMultiIndices"} & Number of included vectors before
resonant grouping and cancellation.\\
\bottomrule
\end{longtable}
The expression is not wrapped in \code{ConditionalExpression}; its domain
and branch are part of the documented result contract. It should be
evaluated only for $y>0$ and parameter values satisfying the assumptions.
All asymptotic assertions concern $y\to0^+$.

For example, the first nontrivial logarithmic output is represented by
\begin{lstlisting}
{{1, 1}, {2, -1 - LogCoordinate}}
\end{lstlisting}
with remainder exponent $3$ and logarithmic degree $2$. For the second
example, \code{4 Sqrt[2] - 3} excludes the fifth block exactly; no
floating-point approximation to that cutoff is required.

\subsection{\texttt{InverseResidual}}
Use
\begin{lstlisting}
InverseResidual[result]
InverseResidual[result, explicitResidualCutoff]
InverseResidual[result, Automatic,
  "MaxCompositionOrder" -> 128]
\end{lstlisting}
The default cutoff is \eqref{eq:residual-cutoff}. The result stores the
coefficient pairs in \code{"ResidualBlocks"}. Its expression fields are
\code{"NormalizedResidualExpression"} and \code{"ResidualExpression"};
the latter is the normalized residual jet multiplied by $a$. It also
returns \code{"Cutoff"} and the boolean \code{"VanishingBelowCutoff"}.

The routine expands only ordinary analytic units in an auxiliary
variable, using \code{Series} where ordinary Taylor series are appropriate.
It does not ask \code{Series} to solve an inverse at the singular endpoint.
This distinction also explains why using an association of exact blocks
is preferable to forcing the answer into \code{SeriesData}. The documented
\code{SeriesData} representation has a single evenly stepped power index
\cite{seriesdata}; the package needs an arbitrary exactly ordered
positive semigroup and coefficient polynomials in a logarithm.

A failed composition identity is informative: a wrong coefficient, an
omitted resonant contribution, or an incorrect logarithmic shift leaves
a nonzero residual block. The test suite deliberately corrupts a stored
coefficient to ensure that the checker is designed to detect such an
error rather than merely returning an expected status.

\subsection{\texttt{LagrangeInverseTruncation}}
The low-level helper is
\begin{lstlisting}
LagrangeInverseTruncation[h, {x, t}, n]
LagrangeInverseTruncation[h, {x, t}, n,
  "OutputPower" -> q]
\end{lstlisting}
Here $h$ is the correction in $x+h(x)$, not the entire forward function.
The helper returns
\begin{equation}\label{eq:helper}
 t^q+\sum_{k=1}^{n}\frac{q(-1)^k}{k!}
 \left.\left(\frac{\dd}{\dd x}\right)^{k-1}
       \bigl(x^{q-1}h(x)^k\bigr)\right|_{x=t}.
\end{equation}
For $q=1$ it is the near-identity inverse truncation. It accepts broader
symbolic expressions and symbolic exponent parameters because it does
not enumerate or order their output support. It does \emph{not} prove
$h=o(x)$, derivative estimates, convergence at the auxiliary parameter
$1$, or an asymptotic remainder. Those hypotheses remain the caller's
responsibility. This separation prevents a general formal identity from
being misrepresented as a universal certified asymptotic solver.

\subsection{Numerical evaluation and use with parameters}
A typical exact-to-numeric workflow is
\begin{lstlisting}
Clear[x, y, a, c];
r = RealInverseExpansion[a x + c x^2, {x, y}, 5,
  Assumptions -> a > 0 && Element[c, Reals]];
N[r["Expression"] /. {a -> 3, c -> -2, y -> 10^-8}, 80]
\end{lstlisting}
Do not start with \lstinline|10.^-8| and later raise its displayed precision:
that does not recover digits that were never present. Likewise, an
approximate exponent can destroy the exact identity of resonant powers.
For numerical comparisons, sufficient guard digits should be used when
subtracting a small residual or an error much smaller than the inverse
itself.

The directory can be loaded directly with \code{Get}; no installation,
network service, or third-party Wolfram add-on is required for execution
in a suitable kernel. An optional \code{PacletInfo.wl} and conventional
\code{Kernel/init.m} are supplied. The declared compatibility target is
Wolfram Language 12.0 or later; this is an implementation target, not a
claim of executed cross-version testing.

\section{Verification and reproducibility}\label{sec:verification}

\subsection{What was actually executed}
The delivered independent verification script ran under Python 3.13.5,
SymPy 1.14.0, and mpmath 1.3.0, with 180 decimal digits for numerical
comparisons. Its recorded result is \textbf{66 checks passed out of 66}.
The result files identify each check individually.

The checks comprise eight logarithmic coefficients tested in three ways
(implicit recursion, direct differentiation, and the Catalan leading-log
identity); eight irrational-power coefficients checked by a separate
implicit recursion; eight independent composition checks for different
models; four support, remainder, resonance, and corruption checks; four
static delimiter checks; and eighteen high-precision numerical comparisons.
The composition models include mixed irrational/logarithmic support,
resonance, nonlinear leading powers, nonunit leading coefficients,
a pure monomial, and quadratic logarithmic coefficients.

The unit-series recurrences used independently are worth recording. If
$A(s)=1+\sum_{n\ge1}a_ns^n$ and $A(s)^\theta=\sum_{n\ge0}b_ns^n$, then
$A B'=\theta B A'$ gives
\begin{equation}\label{eq:unit-power-recurrence}
 b_0=1,\qquad
 nb_n=\sum_{k=1}^{n}\bigl((\theta+1)k-n\bigr)a_kb_{n-k}.
\end{equation}
If $\log A(s)=\sum_{n\ge1}\ell_ns^n$, the identity
$A(\log A)'=A'$ gives
\begin{equation}\label{eq:unit-log-recurrence}
 \ell_n=a_n-\frac1n\sum_{k=1}^{n-1}k\ell_k a_{n-k}.
\end{equation}
These recurrences do not call the inverse coefficient formula. They
provide an algebraically independent route to the unit compositions and
implicit coefficient solutions used in the checks.

\subsection{Numerical results}
Table~\ref{tab:numeric} reports relative errors against independently
computed positive roots. Here $N$ is the number of correction orders, so
the leading term is additional. For the logarithmic example, retaining
$N$ corrections means powers through $y^{N+1}$; for the irrational example,
it means powers through $y^{1+N(\sqrt2-1)}$.
\begin{table}[htbp]
\centering
\small
\begin{tabular}{@{}llrrr@{}}
\toprule
Example & $y$ & $N=2$ & $N=4$ & $N=6$\\
\midrule
Logarithmic & $10^{-2}$ & $1.746\times10^{-4}$ & $1.242\times10^{-6}$ & $1.039\times10^{-8}$ \\
Logarithmic & $10^{-4}$ & $2.408\times10^{-9}$ & $1.150\times10^{-14}$ & $6.630\times10^{-20}$ \\
Logarithmic & $10^{-8}$ & $2.478\times10^{-20}$ & $5.850\times10^{-33}$ & $1.674\times10^{-45}$ \\
Irrational & $10^{-2}$ & $6.806\times10^{-3}$ & $4.740\times10^{-4}$ & $3.852\times10^{-5}$ \\
Irrational & $10^{-4}$ & $2.414\times10^{-5}$ & $3.764\times10^{-8}$ & $6.802\times10^{-11}$ \\
Irrational & $10^{-8}$ & $2.625\times10^{-10}$ & $1.994\times10^{-16}$ & $1.753\times10^{-22}$ \\
\bottomrule
\end{tabular}
\caption{Observed relative errors, computed at 180 decimal digits and rounded for display. The reference roots were computed independently by Newton iteration. These are numerical diagnostics, not interval certificates.}
\label{tab:numeric}
\end{table}


The slower improvement for the irrational example at a fixed $y$ is
expected: its small parameter is $y^{\sqrt2-1}$ rather than approximately
$y|\log y|$. The table is a numerical consistency check, not a substitute
for the error proofs and not a rigorous floating-point enclosure.
Full values, absolute errors, and residual magnitudes are retained in
\code{Verification/results.json}.

\subsection{What was not executed}
The archive supplies \textbf{34 native Wolfram verification tests} and a
\code{wolframscript} runner. Those tests were \textbf{not executed} in the
preparation environment. The available Wolfram service returned HTTP 404,
and no local Wolfram kernel was installed. Static delimiter checks are
not a parser, a type checker, or kernel execution. The Python reference
implementation verifies the mathematics and its own calculations; it
does not establish runtime correctness of the Wolfram source.

This distinction is part of the delivery record, not an omitted
qualification. The package has been written and reviewed against the
specified Wolfram semantics, and the native tests are intended to make
an actual kernel run straightforward and auditable. Until such a run is
recorded, compatibility and runtime behavior should not be described as
empirically verified. No current result of \code{AsymptoticSolve},
\code{InverseSeries}, or the historical \code{InverseFunction} calls is
claimed to have been reproduced.

\subsection{Commands to reproduce the calculations and document}
From the extracted top-level directory:
\begin{lstlisting}[language={}]
python Verification/verify.py
wolframscript -file Tests/run.wls
cd Article
pdflatex -interaction=nonstopmode -halt-on-error real_inverse_asymptotics.tex
pdflatex -interaction=nonstopmode -halt-on-error real_inverse_asymptotics.tex
pdflatex -interaction=nonstopmode -halt-on-error real_inverse_asymptotics.tex
\end{lstlisting}
The Python script needs SymPy and mpmath. Its output records software
versions and rewrites the independent results files. The Wolfram runner
prints the kernel version and test counts and writes a native test report;
no such report is pre-populated in this archive. The article builds with
a standard TeX Live installation containing the packages named in its
preamble. Its source and the numerical-table fragment are included.

\section{Extensions and precise boundaries}\label{sec:limits}

\subsection{What the theory and implementation already cover}
Within the exact finite-model class, all coefficients are covered, not
just the displayed examples. Arbitrary positive real leading powers,
multiple exact irrational corrections, logarithmic polynomials of any
finite degree, symbolic real coefficients under assumptions, and resonant
exponent sums are handled by the same formula. Any desired finite
exclusive power cutoff determines a finite mathematical computation.
The implementation can stop early only with an explicit resource or
undecidability failure, rather than silently claiming the requested order.

The output class is a finite-generator power--logarithmic subclass of
broader transseries calculi; Edgar's exposition gives context for those
larger structures \cite{edgar}. No claim of a complete implementation of
the full logarithmic-exponential transseries field is made.

\subsection{Leading logarithms and other dominant cores}
A function whose leading term contains a logarithm requires another
normalization. For instance,
\[
 y=-x\log x,\qquad x\to0^+,
\]
has the small positive inverse
\[
 x=\exp(W_{-1}(-y))=-\frac{y}{W_{-1}(-y)}.
\]
To derive this, put $u=\log x$ and solve $u e^u=-y$; the branch
$u\to-\infty$ is $W_{-1}$. This inverse brings in inverse logarithms
and iterated logarithms. It is not a finite power--log perturbation of a
pure positive monomial in the sense of \eqref{eq:model}.
The present front end rejects such input instead of choosing the wrong
Lambert branch. A future dominant-core layer could supply the appropriate
exact inverse and transport corrections around it, in the spirit of the
repository's core-inversion architecture \cite{repo}.

\subsection{Symbolic exponents and parameter chambers}
If an exponent is a parameter, support order can change with that
parameter. For example, comparing $2\delta_1$ with $\delta_2$ requires
knowing on which side of $2\delta_1=\delta_2$ the parameters lie.
A fully symbolic ordered solver should decompose parameter space into
chambers and treat the resonance boundaries separately. Merely knowing
that both weights are positive is insufficient to choose one ordered
list of the first few blocks. The derivative helper remains useful in
this situation, but its perturbation grading is not an automatic sorted
asymptotic grading.

\subsection{Infinite supports, inverse logarithms, and flat terms}
Finitely many strictly positive weights are what make every cutoff a
finite task. If infinitely many generators approach zero, a bounded
power interval can contain infinitely many independent contributions.
Different support hypotheses and truncation semantics are then needed.
Similarly, allowing rational or asymptotic functions of $L$ rather than
polynomials can create infinitely many logarithmic orders at a fixed
power. Those extensions require a second, explicit coefficient-scale
truncation, not just removal of the polynomial check in the parser.

A function and its perturbation by a term such as $e^{-1/x}$ may share
all power--log coefficients while having different exact inverses.
Under the derivative conditions of Section~\ref{sec:input-error}, their
inverse difference is correspondingly flat. Thus a formal expansion alone
does not uniquely determine an arbitrary smooth inverse. The convergence
theorem here identifies the inverse because the finite forward model is
specified exactly, rather than only up to a flat error.

\subsection{A possible formal verification boundary}
The finite coefficient engine has a relatively small mathematical trust
boundary: polynomial arithmetic over exact coefficients; the block
operator identity; finite weighted-support enumeration; exact exponent
comparisons; and the finite composition identity. These pieces are
natural candidates for independent formal verification. The analytic
branch, convergence, and remainder theorems form a separate layer that
connects those finite computations to real functions. Formalizing only
the coefficient identities would not by itself certify that analytic
connection, a distinction emphasized in the repository's own status
records \cite{repovol}.

\section{Conclusion}
The right replacement for a Taylor expansion of an unspecified complex
inverse is an expansion of a \emph{specified positive real germ} on its
correct ordered scale. For the two questions at hand, the complete
solution is supplied by local Lagrange inversion around a positive point,
followed by a uniform small-parameter argument. Polynomial logarithmic
coefficients are managed by constant-coefficient differential operators;
irrational powers are managed by a locally finite positive semigroup.
Neither feature requires numerical guessing or a symbolic global inverse.

The resulting package keeps the three logically different tasks separate:
computing exact coefficient blocks, certifying their finite composition
identity, and proving the asymptotic meaning of the omitted tail. That
separation is what makes the original real-branch examples, their mixed
generalizations, and nonunit leading powers accessible to the same
explicit algorithm.

\appendix
\section{Additional logarithmic coefficients}\label{app:coefficients}
The polynomials $A_n$ in \eqref{eq:log-polys} through $n=4$ appear in
\eqref{eq:log-answer}. The next two are
\begin{align*}
 A_5(L)={}&-42L^5-\frac{3727}{12}L^4-\frac{5365}{6}L^3
                -1256L^2-\frac{5177}{6}L-\frac{2789}{12},\\
 A_6(L)={}&132L^6+\frac{5981}{5}L^5+\frac{52937}{12}L^4
                +\frac{16979}{2}L^3\\
          &+9007L^2+\frac{30019}{6}L+\frac{22769}{20}.
\end{align*}
Keeping these through $n=6$ gives error
$\OO(y^8\ell(y)^7)$. Coefficients through $n=8$ are recorded in
\code{Verification/results.json}; the operator formula produces any
further fixed order. The leading coefficients $-42$ and $132$ are
$(-1)^5\operatorname{Cat}_5$ and
$(-1)^6\operatorname{Cat}_6$, respectively.

For numerical evaluation it is often more efficient to evaluate each
coefficient polynomial by Horner's rule than to expand the entire
expression in the numerical variable. The supplied symbolic package
prioritizes exact readable output, while arbitrary-precision evaluation
can be applied after the exact substitution.

\clearpage
\section{Archive organization and trust record}
\begin{lstlisting}[language={}]
RealInverseAsymptotics/
  README.md
  PacletInfo.wl
  Kernel/
    RealInverseAsymptotics.wl
    init.m
  Examples/
    examples.wl
  Tests/
    RealInverseAsymptotics.wlt
    run.wls
  Verification/
    verify.py
    results.json
    results.txt
  Article/
    real_inverse_asymptotics.tex
    real_inverse_asymptotics.pdf
    numerical_table.tex
    build.sh
  MANIFEST.sha256
\end{lstlisting}
The manifest records checksums for delivered files other than the
manifest itself. No external paper, repository volume, font file, or
proprietary kernel is redistributed. The article is self-contained apart
from ordinary LaTeX package dependencies; its proofs do not depend on
having access to the linked repository.

The Wolfram tests cover the two original examples, logarithmic error
metadata, exclusive irrational cutoffs, Catalan and resonant cases,
positive power cores, symbolic coefficients and slopes, residual
composition, deliberate corruption, the general derivative helper,
invalid inputs, resource limits, and one high-precision root comparison.
Their expected outputs are specifications awaiting an actual native
kernel run, not fabricated execution receipts.

\clearpage
\begin{thebibliography}{9}
\bibitem{question}
Vladimir Reshetnikov,
\emph{How to get an asymptotic of the real-valued branch of the inverse
function?}, Mathematica Stack Exchange, question 236367, posted
11 December 2020, with subsequent update.
\url{https://mathematica.stackexchange.com/questions/236367/how-to-get-an-asymptotic-of-the-real-valued-branch-of-the-inverse-function}.
The supplied archive was the primary input; the public page was also
consulted on 7 September 2026.

\bibitem{dlmf}
NIST Digital Library of Mathematical Functions,
\S1.10(vii), \emph{Inverse Functions}, Lagrange Inversion Theorem and
Extended Inversion Theorem.
\url{https://dlmf.nist.gov/1.10#vii}.
Consulted 7 September 2026.

\bibitem{gessel}
Ira M. Gessel, \emph{Lagrange Inversion},
Journal of Combinatorial Theory, Series A \textbf{144} (2016), 212--249.
\url{https://arxiv.org/abs/1609.05988}.
The arXiv version's Theorem 2.1.1 states the coefficient form used here.

\bibitem{repo}
Vladimir Reshetnikov et al., ProveIt repository,
\emph{Series and transseries}, group README.
\url{https://github.com/VladimirReshetnikov/ProveIt/tree/main/Analysis/FabiusFunction/docs/semi-formalized-research-frontiers/drafts/series-and-transseries}.
Consulted 7 September 2026. Context used: comparison of operator and
coefficient inversion, dominant-core architecture, and residual transport.

\bibitem{repovol}
Vladimir Reshetnikov et al., ProveIt repository,
\emph{Transseries: the polynomial--logarithmic calculus and its inversions},
canonical-volume README.
\url{https://github.com/VladimirReshetnikov/ProveIt/tree/main/Analysis/FabiusFunction/docs/semi-formalized-research-frontiers/drafts/series-and-transseries/Transseries_And_Inversion}.
Consulted 7 September 2026. Context used: scoped formalization inventory
and explicit boundaries between formal identities and analytic claims.

\bibitem{inversefunction}
Wolfram Research, \emph{InverseFunction}, Wolfram Language documentation.
\url{https://reference.wolfram.com/language/ref/InverseFunction.html}.
Consulted 7 September 2026.

\bibitem{seriesdata}
Wolfram Research, \emph{SeriesData}, Wolfram Language documentation.
\url{https://reference.wolfram.com/language/ref/SeriesData.html}.
See also \emph{InverseSeries},
\url{https://reference.wolfram.com/language/ref/InverseSeries.html}.
Consulted 7 September 2026.

\bibitem{edgar}
G. A. Edgar, \emph{Transseries for Beginners},
Real Analysis Exchange \textbf{35} (2009/2010), 253--309.
\url{https://arxiv.org/abs/0801.4877}.
The real grid-based viewpoint provides broader context; the finite-model
proofs in this article are given independently.
\end{thebibliography}
\end{document}
